% !TEX root = ../../hw07.tex
% Source: author's supplied text in chat, 2026-11-06, Q3.
% Retain the auxiliary primitive and the endpoint calculation.
\begin{proof}[Solution]
  Since $P(x)$ represents a polynomial with real coefficients,
  \[
    P(x)=a_0+a_1x+\cdots+a_nx^n,
  \]
  it is differentiable and continuous.
  If $P(0)=P(1)$, Rolle's theorem would give a
  point $x\in(0,1)$ with $P'(x)=0$. However, we are interested in
  $P(x)=0$, so we can construct an auxiliary polynomial $F(x)$
  such that $F'(x)=P(x)$.

  Since $F(x)$ is a polynomial, it is continuous on $[0,1]$ and
  differentiable on $(0,1)$. If we can establish that $F(0)=F(1)$,
  Rolle's theorem guarantees some $x\in(0,1)$ satisfying
  $F'(x)=0$ and therefore $P(x)=0$.

  For the derivative of $F(x)$ to be equal to $P(x)$, we can select
  terms by inverting the power rule $(x^n)'=nx^{n-1}$,
  \[
    F(x)=a_0x+\frac{a_1}{2}x^2+\cdots+
      \frac{a_n}{n+1}x^{n+1}+C,\qquad C\in\R.
  \]
  Evaluating $F(x)$ at both endpoints gives
  \[
    \begin{aligned}
      F(0)&=C,\\
      F(1)&=\underbrace{a_0+\frac{a_1}{2}+\cdots+
        \frac{a_n}{n+1}}_{=0}+C.
    \end{aligned}
  \]
  Thus $F(0)=F(1)=C$. By Rolle's theorem, there exists
  $x\in(0,1)$ such that $F'(x)=0$. Since $F'(x)=P(x)$, we conclude
  that $P(x)=0$, as required.
\end{proof}
